\section{A first calculation}
\label{sec:first-calculation}

This chapter computes a CASSCF nuclear gradient for water at the
6-31G level, with an active space of two electrons in two orbitals
(the highest occupied and lowest unoccupied orbital of the reference).
It is the smallest complete run through the kernel: a reference SCF, an
orbital-rotation optimization, and one nuclear gradient from automatic
differentiation.

\subsection{The script}

\begin{lstlisting}[language=Python]
import os
os.environ['JAX_PLATFORMS'] = 'cpu'
for v in ('OMP_NUM_THREADS', 'MKL_NUM_THREADS',
          'OPENBLAS_NUM_THREADS', 'NUMEXPR_NUM_THREADS',
          'VECLIB_MAXIMUM_THREADS'):
    os.environ[v] = '8'

import numpy as np
import jax
jax.config.update('jax_enable_x64', True)
from jax import numpy as jnp
from pyscfad import gto as pgto, scf as pscf

from diffmrc.core.rotation import block_rotation
from diffmrc.core.orblib import orbital_pairs
from diffmrc.core.optimizer import optimize_theta
from diffmrc.core.grad import (make_energy_fn, make_grad_fn,
                               solve_theta_star, nuclear_gradient)

mol = pgto.Mole()
mol.atom = 'O 0 0 0; H 0.7557 0 0.5880; H -0.7557 0 0.5880'
mol.basis = '6-31g'
mol.build()

nmo = mol.nao_nr()
ncore, nact, nele = 4, 2, 2
pairs = orbital_pairs(ncore, nact, nmo, include_cv=True)

E_cas = make_energy_fn(ncore, nact, nele)
mf = pscf.RHF(mol)
mf.kernel()
C_rhf = jnp.asarray(np.asarray(mf.mo_coeff))

h1e = mol.intor('int1e_kin') + mol.intor('int1e_nuc')
eri = mol.intor('int2e', aosym='s1')

E_theta = lambda t: E_cas(C_rhf @ block_rotation(t, nmo, pairs), h1e, eri,
                          mol.energy_nuc())
th_star, gnorm = solve_theta_star(mol, C_rhf, E_theta, len(pairs),
                                  optimize_theta, steps=2000,
                                  g_conv=1e-12)
g = nuclear_gradient(mol, make_grad_fn(nmo, pairs, th_star, E_cas, pscf))
print(g)          # shape (natm, 3), in Eh/bohr
\end{lstlisting}

\subsection{What each piece does}

The reference SCF is an ordinary spin-restricted Hartree--Fock from
pyscfad, whose MO coefficients are the starting point of the chain.
The rotation parameter vector \texttt{t} is turned into an orbital
rotation by \texttt{block\_rotation}; the pair list from
\texttt{orbital\_pairs} includes the core--virtual rotations, which are
required for a correct stationary point --- omitting them leaves the
optimizer stopped at a point that is not stationary.  \texttt{E\_cas}
is the CASCI energy of the rotated orbitals.

\texttt{solve\_\allowbreak theta\_\allowbreak star} runs the two-stage orbital optimization
and returns the parameter vector \emph{and} the gradient norm at the
point it stopped at.  Do not discard \texttt{gnorm}: it is the
evidence of convergence, and the package reports stationarity from
this quantity, not from an optimizer status code.

\texttt{nuclear\_gradient} applies automatic differentiation to the
whole converged chain with respect to the nuclear coordinates, with
the orbital parameters held at their converged values.  The mathematics
of that construction is in the Manual (Chapter on the differentiable
chain); for this tutorial the point is that no derivative expression
was written anywhere in the package or in this script.

\subsection{Reading the result}

For this geometry the gradient is small but not zero --- the atoms are
close to, not at, the CASSCF/6-31G equilibrium.  Two checks belong to
every such run:

\begin{itemize}
  \item \texttt{gnorm}, the orbital-rotation gradient norm, should sit
    at or below the threshold that
    \texttt{solve\_\allowbreak theta\_\allowbreak star} was asked for.  If the refinement
    stalls above the threshold, the returned point carries the stall
    in its stopping reason; treat a stalled point as a candidate, not
    as a result.
  \item The energy is invariant under rotations within a degenerate
    orbital pair.  If a molecule has exact degeneracies, small
    differences between two runs with different rotation orders are a
    property of the chart, not a failure of the code.
\end{itemize}

The next chapter moves from the ground state to several states at once
--- the setting in which the state-averaged machinery and the
per-state perturbation theory come in.
